\documentclass[10pt,a4paper]{amsart}
\usepackage{amsmath,amssymb,amsthm,graphicx}
\usepackage{verbatim}
\newenvironment{reference}{\comment}{\endcomment}
%\newenvironment{reference}{}{}
\newenvironment{slogan}{\comment}{\endcomment}
\newenvironment{history}{\comment}{\endcomment}

% For commutative diagrams we use Xy-pic
\usepackage[all]{xy}

% We use 2cell for 2-commutative diagrams.
\xyoption{2cell}
% Theorem environments.
%
\theoremstyle{plain}
\newtheorem{theorem}[subsection]{Theorem}
\newtheorem{proposition}[subsection]{Proposition}
\newtheorem{lemma}[subsection]{Lemma}

\theoremstyle{definition}
\newtheorem{definition}[subsection]{Definition}
\newtheorem{example}[subsection]{Example}
\newtheorem{exercise}[subsection]{Exercise}
\newtheorem{situation}[subsection]{Situation}

\theoremstyle{remark}
\newtheorem{remark}[subsection]{Remark}
\newtheorem{remarks}[subsection]{Remarks}

\numberwithin{equation}{subsection}

% Macros
%
\def\lim{\mathop{\mathrm{lim}}\nolimits}
\def\colim{\mathop{\mathrm{colim}}\nolimits}
\def\Spec{\mathop{\mathrm{Spec}}}
\def\Hom{\mathop{\mathrm{Hom}}\nolimits}
\def\Ext{\mathop{\mathrm{Ext}}\nolimits}
\def\SheafHom{\mathop{\mathcal{H}\!\mathit{om}}\nolimits}
\def\SheafExt{\mathop{\mathcal{E}\!\mathit{xt}}\nolimits}
\def\Sch{\mathit{Sch}}
\def\Mor{\mathop{\mathrm{Mor}}\nolimits}
\def\Ob{\mathop{\mathrm{Ob}}\nolimits}
\def\Sh{\mathop{\mathit{Sh}}\nolimits}
\def\NL{\mathop{N\!L}\nolimits}
\def\CH{\mathop{\mathrm{CH}}\nolimits}
\def\proetale{{pro\text{-}\acute{e}tale}}
\def\etale{{\acute{e}tale}}
\def\QCoh{\mathit{QCoh}}
\def\Ker{\mathop{\mathrm{Ker}}}
\def\Im{\mathop{\mathrm{Im}}}
\def\Coker{\mathop{\mathrm{Coker}}}
\def\Coim{\mathop{\mathrm{Coim}}}

% Boxtimes
%
\DeclareMathSymbol{\boxtimes}{\mathbin}{AMSa}{"02}

%
% Macros for moduli stacks/spaces
%
\def\QCohstack{\mathcal{QC}\!\mathit{oh}}
\def\Cohstack{\mathcal{C}\!\mathit{oh}}
\def\Spacesstack{\mathcal{S}\!\mathit{paces}}
\def\Quotfunctor{\mathrm{Quot}}
\def\Hilbfunctor{\mathrm{Hilb}}
\def\Curvesstack{\mathcal{C}\!\mathit{urves}}
\def\Polarizedstack{\mathcal{P}\!\mathit{olarized}}
\def\Complexesstack{\mathcal{C}\!\mathit{omplexes}}
% \Pic is the operator that assigns to X its picard group, usage \Pic(X)
% \Picardstack_{X/B} denotes the Picard stack of X over B
% \Picardfunctor_{X/B} denotes the Picard functor of X over B
\def\Pic{\mathop{\mathrm{Pic}}\nolimits}
\def\Picardstack{\mathcal{P}\!\mathit{ic}}
\def\Picardfunctor{\mathrm{Pic}}
\def\Deformationcategory{\mathcal{D}\!\mathit{ef}}
\usepackage{cleveref}
\crefname{originallemma}{original lemma}{original lemmas}
\Crefname{originallemma}{Original lemma}{Original lemmas}
\makeatletter
\newlabel{chow-lemma-K-tensor-Q}{{0FEY}{source}{}{sourcepointer-1}{}}
\newlabel{chow-lemma-K-tensor-Q@cref}{{[originallemma][0][]0FEY}{[1][0][]source}{}{}{}}
\makeatother
\crefname{originallemma}{original lemma}{original lemmas}
\Crefname{originallemma}{Original lemma}{Original lemmas}
\makeatletter
\newlabel{chow-lemma-second-adams-operator}{{0FEJ}{source}{}{sourcepointer-2}{}}
\newlabel{chow-lemma-second-adams-operator@cref}{{[originallemma][0][]0FEJ}{[1][0][]source}{}{}{}}
\makeatother
\crefname{originalremark}{original remark}{original remarks}
\Crefname{originalremark}{Original remark}{Original remarks}
\makeatletter
\newlabel{chow-remark-chern-character-K}{{0FEM}{source}{}{sourcepointer-3}{}}
\newlabel{chow-remark-chern-character-K@cref}{{[originalremark][0][]0FEM}{[1][0][]source}{}{}{}}
\makeatother
\crefname{originalremark}{original remark}{original remarks}
\Crefname{originalremark}{Original remark}{Original remarks}
\makeatletter
\newlabel{chow-remark-gysin-chern-classes}{{0FF8}{source}{}{sourcepointer-4}{}}
\newlabel{chow-remark-gysin-chern-classes@cref}{{[originalremark][0][]0FF8}{[1][0][]source}{}{}{}}
\makeatother
\crefname{originalsection}{original section}{original sections}
\Crefname{originalsection}{Original section}{Original sections}
\makeatletter
\newlabel{chow-section-intersection-regular}{{0FEX}{source}{}{sourcepointer-5}{}}
\newlabel{chow-section-intersection-regular@cref}{{[originalsection][0][]0FEX}{[1][0][]source}{}{}{}}
\makeatother
\crefname{originalreference}{original reference}{original references}
\Crefname{originalreference}{Original reference}{Original references}
\makeatletter
\newlabel{lemma-adams-and-chern}{{0FEN}{source}{}{sourcepointer-6}{}}
\newlabel{lemma-adams-and-chern@cref}{{[originalreference][0][]0FEN}{[1][0][]source}{}{}{}}
\makeatother
\crefname{originalreference}{original reference}{original references}
\Crefname{originalreference}{Original reference}{Original references}
\makeatletter
\newlabel{lemma-chern-classes}{{0FIE}{source}{}{sourcepointer-7}{}}
\newlabel{lemma-chern-classes@cref}{{[originalreference][0][]0FIE}{[1][0][]source}{}{}{}}
\makeatother
\title{The cycle class map from the Chern character on K theory}
\author{The Stacks Project Authors}
\date{stacks-project:weil.tex:local-snapshot-20261012; GNU FDL 1.2 or later}
\begin{document}
\maketitle
\noindent\textit{Editorial note.} All delivered mathematical text is the author's own, copied from the pinned Stacks Project chapter file weil.tex; no statement or proof step was rewritten, reordered or removed.\par\medskip
\noindent\textit{Editorial note.} The counted proof is the author's complete argument: the definition of the cycle class map through the Chern character isomorphism, the pullback compatibility, and the gradedness through the second Adams operator.\par\medskip
\noindent\textit{Editorial note.} The label lemma-cycle-classes names the counted statement, and the proof environment immediately following it in the source is delivered verbatim as its proof.\par\medskip
\noindent\textit{Editorial note.} The counted statement is carried by the author's own lemma environment in the chapter source.\par\medskip
\noindent\textit{Editorial note.} Pre-open gate: python operations/agent-authoring/tools/precard\_gate.py --source weil.tex --label lemma-cycle-classes --start 4712 --end 4758 --local-id 28137 returned RESULT=OK (L28137 free, occupancy FREE, public interval clear). Pre-assembly tag self-check: canonical tag resolved and not present in reference/public-items.jsonl.\par\medskip
\noindent\textit{Editorial note.} The statement and proof are copied from the pinned chapter file; the cited inputs are kept as external pointers with their canonical Stacks tags.\par\medskip
\noindent\textit{Allowed standard prerequisites (editorial).}\begin{itemize}
\item The Chern character identifies K theory with rational Chow groups. Chow Homology, the lemma giving the isomorphism from the Grothendieck group tensored with the rationals to the rational Chow ring via the Chern character capped with the fundamental class, and the remark that it is a ring isomorphism. Provides the isomorphism through which the cycle class map is defined from the Chern character of the cohomology theory.
\item Compatibility of the Chern character with pullbacks. Same chapter, the lemma on the Chern character of the cohomology theory, and Chow Homology, the remark on compatibility of Chern classes with pullbacks. Shows that the constructed cycle class map is compatible with pullbacks.
\item The second Adams operator and the grading of the Chern character. Chow Homology, the lemma on the second Adams operator and its action on classes in a fixed codimension, together with the lemma relating the Adams operators to the Chern character of the cohomology theory. Shows that the cycle class map is graded, i.e. that classes of codimension i land in the cohomological degree 2i.
\end{itemize}\medskip
\ifdefined\hypertarget\hypertarget{sourcepointer-1}{}\fi
\noindent\textit{Source reference.} Original lemma 0FEY: sources/stacks-chow/tags.txt, source lines 19268--19268. Prerequisite pointer invoked at proof line(s) 4731; it names a standard fact of a sibling chapter that the counted argument uses, and every step of the counted proof built on it is retained verbatim in the delivered body.\par\medskip
\ifdefined\hypertarget\hypertarget{sourcepointer-2}{}\fi
\noindent\textit{Source reference.} Original lemma 0FEJ: sources/stacks-chow/tags.txt, source lines 19254--19254. Prerequisite pointer invoked at proof line(s) 4749; it names a standard fact of a sibling chapter that the counted argument uses, and every step of the counted proof built on it is retained verbatim in the delivered body.\par\medskip
\ifdefined\hypertarget\hypertarget{sourcepointer-3}{}\fi
\noindent\textit{Source reference.} Original remark 0FEM: sources/stacks-chow/tags.txt, source lines 19257--19257. Prerequisite pointer invoked at proof line(s) 4732; it names a standard fact of a sibling chapter that the counted argument uses, and every step of the counted proof built on it is retained verbatim in the delivered body.\par\medskip
\ifdefined\hypertarget\hypertarget{sourcepointer-4}{}\fi
\noindent\textit{Source reference.} Original remark 0FF8: sources/stacks-chow/tags.txt, source lines 19278--19278. Prerequisite pointer invoked at proof line(s) 4743; it names a standard fact of a sibling chapter that the counted argument uses, and every step of the counted proof built on it is retained verbatim in the delivered body.\par\medskip
\ifdefined\hypertarget\hypertarget{sourcepointer-5}{}\fi
\noindent\textit{Source reference.} Original section 0FEX: sources/stacks-chow/tags.txt, source lines 19267--19267. Prerequisite pointer invoked at proof line(s) 4754; it names a standard fact of a sibling chapter that the counted argument uses, and every step of the counted proof built on it is retained verbatim in the delivered body.\par\medskip
\ifdefined\hypertarget\hypertarget{sourcepointer-6}{}\fi
\noindent\textit{Source reference.} Original reference 0FEN: sources/stacks-weil/pinned-original-source.tex, source lines 4461--4461. Prerequisite pointer invoked at proof line(s) 4757; it names a standard fact of the same chapter that the counted argument uses, and every step of the counted proof built on it is retained verbatim in the delivered body.\par\medskip
\ifdefined\hypertarget\hypertarget{sourcepointer-7}{}\fi
\noindent\textit{Source reference.} Original reference 0FIE: sources/stacks-weil/pinned-original-source.tex, source lines 4690--4690. Prerequisite pointer invoked at proof line(s) 4734; it names a standard fact of the same chapter that the counted argument uses, and every step of the counted proof built on it is retained verbatim in the delivered body.\par\medskip
\begin{lemma}
\label{lemma-cycle-classes}
Assume given (D0), (D1), and (D2') satisfying axioms (A1), (A2), (A3), and (A4).
There is a unique rule which assigns to every smooth projective $X$ over $k$
a graded ring homomorphism
$$
\gamma : \CH^*(X) \longrightarrow \bigoplus\nolimits_{i \geq 0} H^{2i}(X)(i)
$$
compatible with pullbacks such that $ch^H(\alpha) = \gamma(ch(\alpha))$
for $\alpha$ in $K_0(\textit{Vect}(X))$.
\end{lemma}

\begin{proof}
Recall that we have an isomorphism
$$
K_0(\textit{Vect}(X)) \otimes \mathbf{Q}
\longrightarrow \CH^*(X) \otimes \mathbf{Q},\quad
\alpha \longmapsto ch(\alpha) \cap [X]
$$
see Chow Homology, Lemma \ref{chow-lemma-K-tensor-Q}. It is an isomorphism
of rings by Chow Homology, Remark \ref{chow-remark-chern-character-K}.
We define $\gamma$ by the formula $\gamma(\alpha) = ch^H(\alpha')$
where $ch^H$ is as in Lemma \ref{lemma-chern-classes} and
$\alpha' \in K_0(\textit{Vect}(X))$ is such that
$ch(\alpha') \cap [X] = \alpha$ in $\CH^*(X) \otimes \mathbf{Q}$.

\medskip\noindent
The construction $\alpha \mapsto \gamma(\alpha)$ is compatible
with pullbacks because both $ch^H$ and taking Chern classes
is compatible with pullbacks, see
Lemma \ref{lemma-chern-classes} and
Chow Homology, Remark \ref{chow-remark-gysin-chern-classes}.

\medskip\noindent
We still have to see that $\gamma$ is graded.
Let $\psi^2 : K_0(\textit{Vect}(X)) \to K_0(\textit{Vect}(X))$
be the second Adams operator, see Chow Homology,
Lemma \ref{chow-lemma-second-adams-operator}.
If $\alpha \in \CH^i(X)$ and
$\alpha' \in K_0(\textit{Vect}(X)) \otimes \mathbf{Q}$
is the unique element with $ch(\alpha') \cap [X] = \alpha$,
then we have seen in
Chow Homology, Section \ref{chow-section-intersection-regular}
that $\psi^2(\alpha') = 2^i \alpha'$.
Hence we conclude that $ch^H(\alpha') \in H^{2i}(X)(i)$
by Lemma \ref{lemma-adams-and-chern} as desired.
\end{proof}

\end{document}
